% !TEX program = pdflatex
% ============================================================================
%  Unidade 3 - Customizando o pipeline do spaCy
%  Trabalhando com Texto em Python I  -  CiberExt 26-29 / FEELT38103 / UFU
%
%  COMPILAR COM:  pdflatex unidade3_customizando_pipeline_spacy.tex   (rodar 2x)
%
%  Pacotes necessarios (todos presentes no MacTeX / TeX Live completo):
%    babel, inputenc, fontenc, lmodern, helvet, xcolor, geometry,
%    listings, tcolorbox, titlesec, enumitem, hyperref, microtype
% ============================================================================
\documentclass[11pt,a4paper]{article}

\usepackage[utf8]{inputenc}
\usepackage[T1]{fontenc}
\usepackage[brazilian]{babel}
\usepackage{lmodern}
\usepackage[scaled=0.92]{helvet}
\renewcommand{\familydefault}{\sfdefault}
\usepackage{microtype}
\usepackage[a4paper,margin=2.4cm]{geometry}
\usepackage{xcolor}
\usepackage{enumitem}
\usepackage{listings}
\usepackage[most]{tcolorbox}
\tcbuselibrary{listings,skins,breakable}
\usepackage{titlesec}
\usepackage{hyperref}

% -------------------------------------------------------------- paleta ------
\definecolor{accent}{HTML}{6C4CF1}
\definecolor{cyan}{HTML}{00B4D8}
\definecolor{subink}{HTML}{2B3242}
\definecolor{codebg}{HTML}{1E2430}
\definecolor{codetext}{HTML}{E6E9EF}
\definecolor{ccomment}{HTML}{7FD18B}
\definecolor{cstring}{HTML}{F0B866}
\definecolor{ckw}{HTML}{C792EA}
\definecolor{cbuiltin}{HTML}{82AAFF}
\definecolor{outbg}{HTML}{F0F2F7}
\definecolor{outtext}{HTML}{2F3646}
\definecolor{outframe}{HTML}{9AA4B8}
\definecolor{notec}{HTML}{3B82F6}
\definecolor{tipc}{HTML}{10B981}
\definecolor{warnc}{HTML}{F59E0B}
\definecolor{inlink}{HTML}{5A34D6}

\hypersetup{colorlinks=true,linkcolor=accent,urlcolor=cyan,citecolor=accent,
  pdftitle={Customizando o pipeline do spaCy},pdfauthor={CiberExt 26-29 / FEELT38103 / UFU}}

% --------------------------------------------------- estilo do codigo -------
\lstset{
  extendedchars=true,
  breaklines=true,
  breakatwhitespace=false,
  keepspaces=true,
  columns=fullflexible,
  tabsize=4,
  showstringspaces=false,
  literate=%
    {á}{{\'a}}1 {é}{{\'e}}1 {í}{{\'i}}1 {ó}{{\'o}}1 {ú}{{\'u}}1
    {â}{{\^a}}1 {ê}{{\^e}}1 {ô}{{\^o}}1
    {ã}{{\~a}}1 {õ}{{\~o}}1 {à}{{\`a}}1 {ü}{{\"u}}1 {ä}{{\"a}}1 {ö}{{\"o}}1
    {ç}{{\c c}}1 {Ç}{{\c C}}1
    {Á}{{\'A}}1 {É}{{\'E}}1 {Í}{{\'I}}1 {Ó}{{\'O}}1 {Ú}{{\'U}}1
    {Ã}{{\~A}}1 {Õ}{{\~O}}1 {Â}{{\^A}}1 {Ê}{{\^E}}1 {À}{{\`A}}1
    {—}{{---}}1 {–}{{--}}1
    {“}{{``}}1 {”}{{''}}1 {‘}{{`}}1 {’}{{'}}1 {´}{{'}}1
    {✔}{{[OK]}}4
}

\lstdefinestyle{pystyle}{
  language=Python,
  basicstyle=\ttfamily\small\color{codetext},
  keywordstyle=\color{ckw}\bfseries,
  keywordstyle=[2]\color{cbuiltin},
  morekeywords=[2]{print,type,load,list,append,items,get,spacy,displacy,assert},
  commentstyle=\color{ccomment}\itshape,
  stringstyle=\color{cstring},
}
\lstdefinestyle{outstyle}{
  basicstyle=\ttfamily\small\color{outtext},
}

% caixa de codigo Python (fundo escuro, barra lateral colorida + titulo)
\newtcblisting{pybox}{
  breakable, enhanced,
  listing only, listing options={style=pystyle},
  colback=codebg, colframe=codebg, boxrule=0pt, arc=3pt,
  borderline west={3pt}{0pt}{accent},
  left=10pt,right=10pt,top=6pt,bottom=6pt,
  title=\textbf{Python}, coltitle=white, colbacktitle=accent,
  fonttitle=\sffamily\bfseries\scriptsize,
  before skip=10pt, after skip=12pt,
}
% caixa de saida (fundo claro)
\newtcblisting{outbox}{
  breakable, enhanced,
  listing only, listing options={style=outstyle},
  colback=outbg, colframe=outframe, boxrule=0.4pt, arc=3pt,
  left=10pt,right=10pt,top=6pt,bottom=6pt,
  title=\textbf{Saída}, coltitle=white, colbacktitle=outframe,
  fonttitle=\sffamily\bfseries\scriptsize,
  before skip=10pt, after skip=12pt,
}

% ------------------------------------------------------ caixas de aviso -----
\newtcolorbox{notebox}[1][Nota]{enhanced,breakable,
  colback=notec!6,colframe=notec,coltitle=white,title={#1},
  fonttitle=\sffamily\bfseries\small,boxrule=0pt,leftrule=4pt,arc=3pt,
  left=10pt,right=10pt,top=7pt,bottom=7pt,before skip=10pt,after skip=12pt}
\newtcolorbox{tipbox}[1][Dica]{enhanced,breakable,
  colback=tipc!7,colframe=tipc,coltitle=white,title={#1},
  fonttitle=\sffamily\bfseries\small,boxrule=0pt,leftrule=4pt,arc=3pt,
  left=10pt,right=10pt,top=7pt,bottom=7pt,before skip=10pt,after skip=12pt}
\newtcolorbox{warnbox}[1][Atenção]{enhanced,breakable,
  colback=warnc!8,colframe=warnc,coltitle=white,title={#1},
  fonttitle=\sffamily\bfseries\small,boxrule=0pt,leftrule=4pt,arc=3pt,
  left=10pt,right=10pt,top=7pt,bottom=7pt,before skip=10pt,after skip=12pt}
\newtcolorbox{objbox}[1][Objetivos de aprendizagem]{enhanced,breakable,
  colback=accent!5,colframe=accent,coltitle=white,title={#1},
  fonttitle=\sffamily\bfseries\small,boxrule=0pt,leftrule=4pt,arc=3pt,
  left=10pt,right=10pt,top=7pt,bottom=8pt,before skip=10pt,after skip=14pt}

% ------------------------------------------------------ titulos de secao ----
\titleformat{\section}{\Large\bfseries\sffamily\color{accent}}
  {\thesection.}{0.5em}{}
\titleformat{\subsection}{\large\bfseries\sffamily\color{subink}}
  {\thesubsection}{0.5em}{}
\titlespacing*{\section}{0pt}{18pt}{6pt}
\titlespacing*{\subsection}{0pt}{12pt}{4pt}

% codigo em linha (verbatim, colorido) via atalho @...@
\lstMakeShortInline[basicstyle=\ttfamily\color{inlink}]{@}

\setlength{\parindent}{0pt}
\setlength{\parskip}{0.55em}

\begin{document}

% --------------------------------------------------------------- capa -------
\begin{tcolorbox}[enhanced, colback=accent, coltext=white,
  boxrule=0pt, arc=8pt, left=20pt,right=20pt,top=16pt,bottom=18pt,
  before skip=0pt, after skip=20pt]
  {\footnotesize\bfseries CIBEREXT 26-29 \ \textbullet\ \ FEELT38103 \ \textbullet\ \ UNIVERSIDADE FEDERAL DE UBERLÂNDIA}\\[10pt]
  {\Huge\bfseries Customizando o pipeline do spaCy}\\[8pt]
  {\large Trabalhando com Texto em Python I \ --- \ Unidade 3}\\[6pt]
  {\normalsize Modificar o pipeline, processar em lote, atributos personalizados, salvar em disco e mesclar sintagmas e entidades.}\\[10pt]
  {\footnotesize Adaptado de \textit{Applied Language Technology} (Universidade de Helsinque) --- applied-language-technology.mooc.fi}
\end{tcolorbox}

\begin{objbox}
Ao final desta unidade, você deverá saber:
\begin{itemize}[leftmargin=1.3em,itemsep=2pt,topsep=3pt]
  \item examinar e modificar o pipeline do spaCy;
  \item processar textos com eficiência;
  \item adicionar atributos personalizados a objetos do spaCy;
  \item salvar em disco os textos processados;
  \item mesclar sintagmas nominais e entidades nomeadas em tokens únicos.
\end{itemize}
\end{objbox}

Vamos começar importando a biblioteca spaCy e o módulo @displacy@, usado para desenhar árvores de dependência, e carregando um modelo de linguagem para o inglês:

\begin{pybox}
# Importa a biblioteca spaCy e o modulo displacy
from spacy import displacy
import spacy

# Carrega um modelo pequeno de lingua inglesa e o atribui a variavel 'nlp'
nlp = spacy.load('en_core_web_sm')

# Chama a variavel para examinar o objeto
nlp
\end{pybox}
\begin{outbox}
<spacy.lang.en.English at 0x294522f10>
\end{outbox}

% =========================================================== 1 ==============
\section{Modificando pipelines do spaCy}

O objeto @Language@ é, essencialmente, um \textbf{pipeline} que aplica um modelo de linguagem ao texto, executando as tarefas para as quais o modelo foi treinado. As tarefas realizadas dependem dos \textbf{componentes} presentes no pipeline.

Podemos examinar os componentes com o atributo @pipeline@ de um objeto @Language@:

\begin{pybox}
nlp.pipeline
\end{pybox}
\begin{outbox}
[('tok2vec', <spacy.pipeline.tok2vec.Tok2Vec at 0x294cf3dc0>),
 ('tagger', <spacy.pipeline.tagger.Tagger at 0x294cf3ca0>),
 ('parser', <spacy.pipeline.dep_parser.DependencyParser at 0x294b566d0>),
 ('attribute_ruler', <spacy.pipeline.attributeruler.AttributeRuler at 0x295187680>),
 ('lemmatizer', <spacy.lang.en.lemmatizer.EnglishLemmatizer at 0x295191540>),
 ('ner', <spacy.pipeline.ner.EntityRecognizer at 0x294b56740>)]
\end{outbox}

Isso retorna um objeto @SimpleFrozenList@ do spaCy, composto por \textbf{tuplas} com dois itens: o \textbf{nome} do componente (ex.: @tagger@) e o \textbf{componente} que executa a tarefa (ex.: @spacy.pipeline.tok2vec.Tok2Vec@).

Componentes como @tagger@, @parser@, @ner@ e @lemmatizer@ já são familiares da unidade anterior. Há dois novos:
\begin{itemize}[leftmargin=1.3em,itemsep=2pt,topsep=3pt]
  \item \textbf{\texttt{tok2vec}} mapeia os tokens para suas representações numéricas (veremos na Parte III).
  \item \textbf{\texttt{attribute\_ruler}} aplica regras definidas pelo usuário aos tokens e adiciona essa informação como atributo, se solicitado.
\end{itemize}

\begin{notebox}
A lista em @nlp.pipeline@ \textbf{não} inclui o @Tokenizer@, porque todo texto precisa ser tokenizado para qualquer processamento ocorrer. Por isso ele fica no atributo @tokenizer@, e não em @pipeline@.
\end{notebox}

\subsection{Excluindo componentes para ganhar desempenho}
\textbf{Todo componente tem um custo computacional.} Se você não precisa da saída de um componente, não deve incluí-lo, pois o processamento fica mais lento. Para excluir, forneça o argumento @exclude@ com os nomes dos componentes ao chamar @load()@:

\begin{pybox}
# Carrega um modelo pequeno de lingua inglesa, mas exclui o reconhecimento de
# entidades nomeadas ('ner') e a analise de dependencias sintaticas ('parser').
nlp = spacy.load('en_core_web_sm', exclude=['ner', 'parser'])

# Examina os componentes ativos no objeto Language 'nlp'
nlp.pipeline
\end{pybox}
\begin{outbox}
[('tok2vec', <spacy.pipeline.tok2vec.Tok2Vec at 0x29514d760>),
 ('tagger', <spacy.pipeline.tagger.Tagger at 0x296c7ba00>),
 ('attribute_ruler', <spacy.pipeline.attributeruler.AttributeRuler at 0x296c6bb40>),
 ('lemmatizer', <spacy.lang.en.lemmatizer.EnglishLemmatizer at 0x296c72200>)]
\end{outbox}

Como mostra a saída, @ner@ e @parser@ não estão mais no pipeline.

\subsection{Analisando o pipeline com \textnormal{\texttt{analyze\_pipes()}}}
O método @analyze_pipes()@ dá uma visão geral dos componentes. Definindo @pretty@ como @True@, o spaCy imprime uma tabela com os componentes e as anotações que produzem:

\begin{pybox}
# Analisa o pipeline e guarda a analise na variavel 'pipe_analysis'
pipe_analysis = nlp.analyze_pipes(pretty=True)
\end{pybox}
\begin{outbox}
============================= Pipeline Overview =============================

#   Component         Assigns       Requires   Scores      Retokenizes
-   ---------------   -----------   --------   ---------   -----------
0   tok2vec           doc.tensor                           False
1   tagger            token.tag                tag_acc     False
2   attribute_ruler                                        False
3   lemmatizer        token.lemma              lemma_acc   False

✔ No problems found.
\end{outbox}

O método retorna um \textbf{dicionário} Python. Você pode usá-lo para verificar que nenhum problema foi encontrado \textbf{antes} de processar grandes volumes. Os relatórios ficam sob a chave @problems@:

\begin{pybox}
# Examina o valor guardado sob a chave 'problems'
pipe_analysis['problems']
\end{pybox}
\begin{outbox}
{'tok2vec': [], 'tagger': [], 'attribute_ruler': [], 'lemmatizer': []}
\end{outbox}

Podemos escrever um trecho que \textbf{verifica} isso. Percorremos o dicionário com @items()@ e usamos @assert@ com @len()@ e o operador @==@ para checar que a lista tem tamanho @0@; caso contrário, o Python levanta um @AssertionError@:

\begin{pybox}
# Percorre os pares chave/valor do dicionario. Atribui a chave e o valor
# as variaveis 'component_name' e 'problem_list'.
for component_name, problem_list in pipe_analysis['problems'].items():

    # Usa 'assert' para checar a lista de problemas; levanta Error se necessario.
    assert len(problem_list) == 0, f"There is a problem with {component_name}: {problem_list}!"
\end{pybox}

Usamos uma \textit{f-string}: o @f@ antes das aspas permite \textbf{inserir variáveis} entre chaves @{}@. Se nenhum problema for encontrado, o laço passa silenciosamente.

% =========================================================== 2 ==============
\section{Processando textos com eficiência}

Ao trabalhar com grandes volumes, é desejável processá-los da forma mais eficiente possível. Vamos definir um exemplo com uma \textbf{lista} de três frases da Wikipédia em inglês:

\begin{pybox}
# Reinicializa o modelo de linguagem, pois precisamos da analise de
# dependencias nas proximas secoes.
nlp = spacy.load('en_core_web_sm')

# Define uma lista de frases de exemplo
sents = ["On October 1, 2009, the Obama administration went ahead with a Bush administration program, increasing nuclear weapons production.",
         "The 'Complex Modernization' initiative expanded two existing nuclear sites to produce new bomb parts.",
         "The administration built new plutonium pits at the Los Alamos lab in New Mexico and expanded enriched uranium processing at the Y-12 facility in Oak Ridge, Tennessee."]

# Chama a variavel para examinar a saida
sents
\end{pybox}

Os objetos @Language@ têm um método específico, @pipe()@, para processar textos em uma \textbf{lista}. Ele foi otimizado para processar em \textbf{lotes} (\textit{batches}), sendo mais rápido que um laço @for@. O @pipe()@ recebe uma lista e retorna um \textbf{gerador} (\textit{generator}):

\begin{pybox}
# Alimenta a lista de frases ao metodo pipe()
docs = nlp.pipe(sents)

# Chama a variavel para examinar a saida
docs
\end{pybox}
\begin{outbox}
<generator object Language.pipe at 0x296d0b580>
\end{outbox}

Geradores contêm outros objetos e os \textbf{produzem} (\textit{yield}) quando chamados. Para recuperar todos, convertemos o gerador em uma \textbf{lista}:

\begin{pybox}
# Converte o gerador pipe em uma lista
docs = list(docs)

# Chama a variavel para examinar a saida
docs
\end{pybox}

Isso nos dá uma lista de objetos @Doc@ prontos para processamento posterior.

% =========================================================== 3 ==============
\section{Adicionando atributos personalizados aos objetos do spaCy}

O spaCy permite \textbf{definir atributos personalizados} para objetos @Doc@, @Span@ e @Token@ --- úteis para guardar informações adicionais sobre os textos. Usamos o método @set_extension()@. Como o atributo é adicionado a \textbf{todos} os @Doc@, primeiro importamos o objeto @Doc@ genérico:

\begin{pybox}
# Importa o objeto Doc do modulo 'tokens' do spaCy
from spacy.tokens import Doc

# Adiciona dois atributos personalizados ao objeto Doc, 'age' e 'location',
# usando o metodo set_extension().
Doc.set_extension("age", default=None)
Doc.set_extension("location", default=None)
\end{pybox}

O argumento @default@ define um valor padrão, aqui a palavra-chave @None@.

\begin{notebox}
Diferentemente de atributos como @sents@ ou @heads@, os personalizados ficam sob um atributo que consiste no sublinhado @_@ --- por exemplo, @Doc._.age@.
\end{notebox}

Vamos definir um dicionário aninhado @sents_dict@ com três chaves (@0@, @1@, @2@), cujos valores são dicionários com @age@, @location@ e @text@:

\begin{pybox}
# Cria um dicionario cujos valores sao outros dicionarios
# com tres chaves: 'age', 'location' e 'text'.
sents_dict = {0: {"age": 23,
                  "location": "Helsinki",
                  "text": "The Senate Square is by far the most important landmark in Helsinki."
                 },
              1: {"age": 35,
                  "location": "Tallinn",
                  "text": "The Old Town, for sure."
                 },
              2: {"age": 58,
                  "location": "Stockholm",
                  "text": "Södermalm is interesting!"
                 }
             }
\end{pybox}

Percorremos o dicionário, processamos cada texto e definimos os atributos personalizados:

\begin{pybox}
# Prepara uma lista vazia para guardar os textos processados
docs = []

# Percorre os pares de chave e valor do dicionario 'sents_dict'.
# Os pares chave/valor ficam disponiveis pelo metodo items().
# Chamamos essas chaves e valores de 'key' e 'data'; ou seja, usamos
# a variavel 'data' para se referir ao dicionario aninhado.
for key, data in sents_dict.items():

    # Recupera o valor da chave 'text' do dicionario aninhado.
    # Alimenta esse texto ao modelo em 'nlp' e guarda o resultado em 'doc'.
    doc = nlp(data['text'])

    # Recupera os valores de 'age' e 'location' do dicionario aninhado.
    # Atribui esses valores aos atributos personalizados do objeto Doc.
    # Lembre-se: atributos personalizados ficam sob o pseudo-atributo '_'!
    doc._.age = data['age']
    doc._.location = data['location']

    # Acrescenta o Doc atual em 'doc' a lista 'docs'
    docs.append(doc)
\end{pybox}

Isso produz uma lista de @Doc@s. Vamos percorrê-la e imprimir cada um com seus atributos:

\begin{pybox}
# Percorre cada objeto Doc na lista 'docs'
for doc in docs:

    # Imprime cada Doc e os atributos 'age' e 'location'
    print(doc, doc._.age, doc._.location)
\end{pybox}
\begin{outbox}
The Senate Square is by far the most important landmark in Helsinki. 23 Helsinki
The Old Town, for sure. 35 Tallinn
Södermalm is interesting! 58 Stockholm
\end{outbox}

\subsection{Filtrando dados com compreensão de lista}
Os atributos personalizados podem ser usados para \textbf{filtrar} dados. Uma forma eficiente é a \textbf{compreensão de lista} (\textit{list comprehension}) --- como um laço @for@ declarado ``na hora'', entre colchetes @[]@, com três partes: (1) o que guardar na nova lista; (2) o @for ... in@ percorrendo os itens; (3) o @if@ com a condição:

\begin{pybox}
# Usa uma compreensao de lista para filtrar os Docs cujo atributo
# 'age' tenha valor abaixo de 40.
under_forty = [doc for doc in docs if doc._.get('age') < 40]

# Chama a variavel para examinar a saida
under_forty
\end{pybox}
\begin{outbox}
[The Senate Square is by far the most important landmark in Helsinki.,
 The Old Town, for sure.]
\end{outbox}

Isso retorna uma lista com apenas dois @Doc@s que atendem ao critério.

% =========================================================== 4 ==============
\section{Gravando textos processados em disco}

Com grandes volumes, primeiro garanta que o pipeline funciona com \textbf{poucos} textos; depois processe tudo e \textbf{salve}. O spaCy oferece o tipo @DocBin@ para armazenar objetos @Doc@ com suas anotações:

\begin{pybox}
# Importa o objeto DocBin do modulo 'tokens' do spaCy
from spacy.tokens import DocBin

# Inicializa um objeto DocBin e adiciona os Docs de 'docs'
docbin = DocBin(docs=docs)
\end{pybox}

\begin{warnbox}
Se você adicionou \textbf{atributos personalizados} a @Doc@s, @Span@s ou @Token@s, também é preciso definir @store_user_data=True@ --- por exemplo, @DocBin(docs=docs, store_user_data=True)@.
\end{warnbox}

Verificamos que os três @Doc@s entraram no @DocBin@ com o método @__len__()@:

\begin{pybox}
# Obtem o numero de Docs no DocBin
docbin.__len__()
\end{pybox}
\begin{outbox}
3
\end{outbox}

O método @add()@ permite adicionar mais @Doc@s:

\begin{pybox}
# Define uma string, alimenta o modelo em 'nlp' e adiciona o
# Doc resultante ao objeto DocBin 'docbin'
docbin.add(nlp("Yet another Doc object."))

# Verifica que o Doc foi adicionado; o tamanho agora deve ser 4
docbin.__len__()
\end{pybox}
\begin{outbox}
4
\end{outbox}

Gravamos o @DocBin@ em disco com @to_disk()@, que recebe o argumento @path@:

\begin{pybox}
# Grava o objeto DocBin em disco
docbin.to_disk(path='data/docbin.spacy')
\end{pybox}

Para \textbf{carregar}, inicialize um @DocBin@ vazio e use @from_disk()@:

\begin{pybox}
# Inicializa um novo DocBin e usa 'from_disk' para carregar os dados do disco.
# Atribui o resultado a variavel 'docbin_loaded'.
docbin_loaded = DocBin().from_disk(path='data/docbin.spacy')

# Chama a variavel para examinar a saida
docbin_loaded
\end{pybox}
\begin{outbox}
<spacy.tokens._serialize.DocBin at 0x2953a4c40>
\end{outbox}

Para acessar os @Doc@s guardados, use @get_docs()@, passando o \textbf{vocabulário} (@nlp.vocab@) necessário para reconstruir os dados:

\begin{pybox}
# Usa 'get_docs' para recuperar os Docs do DocBin, passando o
# vocabulario em 'nlp.vocab' para reconstruir os dados.
# Converte o gerador resultante em uma lista para exame.
docs_loaded = list(docbin_loaded.get_docs(nlp.vocab))

# Chama a variavel para examinar a saida
docs_loaded
\end{pybox}
\begin{outbox}
[The Senate Square is by far the most important landmark in Helsinki.,
 The Old Town, for sure.,
 Södermalm is interesting!,
 Yet another Doc object.]
\end{outbox}

\begin{tipbox}[Resumindo]
O ideal é processar os textos \textbf{uma vez}, gravá-los em disco e carregá-los para as análises seguintes.
\end{tipbox}

% =========================================================== 5 ==============
\section{Simplificando a saída: sintagmas nominais e entidades nomeadas}

\subsection{Mesclando sintagmas nominais}
Às vezes é mais vantajoso operar com unidades maiores que tokens individuais, como \textbf{sintagmas nominais} (\textit{noun phrases}). O spaCy dá acesso a eles pelo atributo @noun_chunks@ de um @Doc@:

\begin{pybox}
# Primeiro laco: percorre a lista 'docs'; 'doc' se refere aos itens da lista
for doc in docs:

    # Percorre cada sintagma nominal do objeto Doc
    for noun_chunk in doc.noun_chunks:

        # Imprime o sintagma nominal
        print(noun_chunk)
\end{pybox}
\begin{outbox}
The Senate Square
the most important landmark
Helsinki
The Old Town
Södermalm
\end{outbox}

Para \textbf{mesclar} os sintagmas em um único token, o spaCy oferece a função @merge_noun_chunks@, adicionada ao pipeline com @add_pipe@:

\begin{pybox}
# Adiciona o componente que mescla sintagmas nominais em Tokens unicos
nlp.add_pipe('merge_noun_chunks')
\end{pybox}

\begin{notebox}
Não é preciso reatribuir o objeto @Language@ à mesma variável --- o @add_pipe@ adiciona o componente automaticamente.
\end{notebox}

Processando as frases novamente com @pipe()@, ao percorrer os tokens do primeiro @Doc@ (@[0]@) vemos os sintagmas \textbf{mesclados} e rotulados como @NOUN@:

\begin{pybox}
# Aplica o objeto Language 'nlp' a lista de frases em 'sents'
docs = list(nlp.pipe(sents))

# Percorre os Tokens do primeiro Doc da lista
for token in docs[0]:

    # Imprime o Token e sua classe gramatical
    print(token, token.pos_)
\end{pybox}
\begin{outbox}
On ADP
October PROPN
1 NUM
, PUNCT
2009 NUM
, PUNCT
the Obama administration NOUN
went VERB
ahead ADV
with ADP
a Bush administration program NOUN
, PUNCT
increasing VERB
nuclear weapons production NOUN
. PUNCT
\end{outbox}

Rotular como @NOUN@ é uma aproximação razoável, já que as palavras-núcleo são substantivos. Mesclar os sintagmas \textbf{simplifica} a árvore sintática:

\begin{pybox}
displacy.render(docs[0], style='dep')
\end{pybox}

Os sintagmas continuam disponíveis em @noun_chunks@, guardados como @Span@s com atributos @start@ e @end@:

\begin{pybox}
# Percorre os sintagmas nominais do primeiro Doc [0] na lista 'docs'
for noun_chunk in docs[0].noun_chunks:

    # Imprime o sintagma, seu tipo, e os indices onde comeca e termina
    print(noun_chunk, type(noun_chunk), noun_chunk.start, noun_chunk.end)
\end{pybox}
\begin{outbox}
October <class 'spacy.tokens.span.Span'> 1 2
the Obama administration <class 'spacy.tokens.span.Span'> 6 7
a Bush administration program <class 'spacy.tokens.span.Span'> 10 11
nuclear weapons production <class 'spacy.tokens.span.Span'> 13 14
\end{outbox}

\subsection{Mesclando entidades nomeadas}
Entidades nomeadas são mescladas do mesmo modo, com @merge_entities@. Primeiro removemos a função @merge_noun_chunks@ com @remove_pipe()@:

\begin{pybox}
# Remove a funcao 'merge_noun_chunks' do pipeline em 'nlp'
nlp.remove_pipe('merge_noun_chunks')

# Processa as frases originais novamente
docs = list(nlp.pipe(sents))
\end{pybox}

Em seguida, adicionamos @merge_entities@ e reprocessamos, examinando o terceiro @Doc@:

\begin{pybox}
# Adiciona a funcao 'merge_entities' ao pipeline
nlp.add_pipe('merge_entities')

# Processa os dados novamente
docs = list(nlp.pipe(sents))

# Percorre os Tokens do terceiro Doc da lista
for token in docs[2]:

    # Imprime o Token e sua classe gramatical
    print(token, token.pos_)
\end{pybox}
\begin{outbox}
The DET
administration NOUN
built VERB
new ADJ
plutonium NOUN
pits NOUN
at ADP
the DET
Los Alamos PROPN
lab NOUN
in ADP
New Mexico PROPN
and CCONJ
expanded VERB
enriched ADJ
uranium NOUN
processing NOUN
at ADP
the DET
Y-12 NUM
facility NOUN
in ADP
Oak Ridge PROPN
, PUNCT
Tennessee PROPN
. PUNCT
\end{outbox}

Entidades nomeadas com vários tokens --- como os topônimos \textit{``Los Alamos''} e \textit{``New Mexico''} --- foram \textbf{mescladas} em tokens únicos.

% ======================================================= resumo ============
\section*{Resumo da unidade}
\addcontentsline{toc}{section}{Resumo da unidade}
\begin{objbox}[Nesta unidade, você aprendeu a]
\begin{enumerate}[leftmargin=1.4em,itemsep=2pt,topsep=3pt]
  \item examinar o pipeline com @nlp.pipeline@ e entender @tok2vec@/@attribute_ruler@;
  \item excluir componentes (@exclude=[...]@) e auditar o pipeline com @analyze_pipes()@ e @assert@;
  \item processar em lote com @nlp.pipe()@, entendendo geradores e conversão para lista;
  \item adicionar atributos personalizados (@set_extension()@, @Doc._.@) e filtrar com compreensão de lista;
  \item gravar e carregar com @DocBin@ (@to_disk()@, @from_disk()@, @get_docs()@);
  \item mesclar sintagmas nominais e entidades com @add_pipe()@/@remove_pipe()@.
\end{enumerate}
\end{objbox}

\begin{notebox}[Exercícios sugeridos]
\begin{enumerate}[leftmargin=1.4em,itemsep=3pt,topsep=3pt]
  \item Carregue o modelo excluindo apenas o @lemmatizer@ e confirme, com @nlp.pipeline@, que ele não está mais presente.
  \item Adicione um atributo personalizado @source@ ao @Doc@ e preencha-o com o nome do arquivo de origem ao processar textos da Unidade 1.
  \item Processe as notícias da Unidade 1 com @nlp.pipe()@, grave em um @DocBin@ e recarregue do disco, conferindo o número de @Doc@s.
  \item Compare a árvore de dependências de uma frase \textbf{com} e \textbf{sem} @merge_noun_chunks@ e descreva a diferença.
\end{enumerate}
\end{notebox}

\vspace{6pt}
\noindent\rule{\linewidth}{0.4pt}\\[4pt]
{\footnotesize\textbf{Referências.} \textit{Applied Language Technology} (MOOC), Universidade de Helsinque --- \href{https://applied-language-technology.mooc.fi/}{applied-language-technology.mooc.fi}. Honnibal, M.; Montani, I. et al. \textit{spaCy: Industrial-strength Natural Language Processing in Python} --- \href{https://spacy.io/}{spacy.io}.\par
Material produzido para o Programa CiberExt 26-29 / FEELT38103 --- Universidade Federal de Uberlândia.}

\end{document}
